\scanpage{315}
\unsure{ensure} the local existence of $\tau$, it remains to make sure that
$\varepsilon\sigma$ is a regular section in the sense of Steinberg i.e. is nowhere
a scalar, which follows from the fact that this is so for $\sigma$.

Thus, under the present conditions, let $N'_\sigma\subset N_\sigma$ be the subscheme of
\deleted{$N_\sigma$} $G$ (or of $\unsure{N_\sigma|N}$) defined by
\[
\tag{22}
N'_\sigma=\{\tau\mid[\tau,\sigma]\in\Gamma\mu_2\}
\]
then $N'_\sigma$ is stable under right translation by $T_\sigma$ --- better,
it is a subgroup scheme of $N_\sigma$ \inserted{(containing $T_\sigma$)}, for if
$\tau,\tau'$ are such that
\[
\tau(\sigma)=\varepsilon\sigma,\qquad\tau'\sigma=\varepsilon'\sigma,
\]
we will have
\[
\tau'\tau(\sigma)=\varepsilon'\varepsilon\cdot\sigma,
\]
at the same time we see that
\[
\tau\longmapsto\varepsilon
\]
is a group homomorphism $N'_\sigma\longrightarrow\mu_2$, and we will have an exact sequence of
group schemes
\[
\tag{23}
1\longrightarrow T_\sigma\hookrightarrow N'_\sigma
\xrightarrow{\varepsilon_\sigma}\mu_2\longrightarrow1.
\]
We \unsure{will note} that if $\tau\in\Gamma G$ normalizes $T_\sigma$, we will have
\[
\tau(\sigma)=a\sigma+b\qquad a,b\in\Gamma\mathcal O_S,
\]

\scanpage{316}
\mspage{\unsure{458|488}}
and taking the traces of the two sides we find (taking into account that
$\operatorname{Tr}\tau(\sigma)=\operatorname{Tr}\sigma=0$)
$\operatorname{Tr}b.1=0$ i.e. $2b=0$, so if $b$ is invertible
$b=0$, hence $\tau(\sigma)=a\sigma$ and we necessarily have $a^2=1$ i.e. $a$ is an
$\varepsilon$. Thus
\[
\tag{24}
N'_\sigma=N_\sigma\quad\text{in residual characteristic }\ne2.
\]

\deleted{\unsure{But in characteristic}}
More \unsure{generally}, over a base $S$ such that
$2$ \deleted{not} be \deleted{\unsure{a zero divisor}} \inserted{a regular section}
(i.e. not a zero divisor) in $\Gamma(S,\mathcal O_S)$, we see that every section of
$\unsure{N'_\sigma|N_\sigma}$ is in fact a section of $N_\sigma$.
But this is no longer so in characteristic $2$, for
\[
\sigma'=a\sigma+b\qquad(a\in\Gamma G_m,\ b\in\Gamma\mathcal O_S)
\]
as soon as it satisfies
\[
\operatorname{Tr}\sigma'=\unsure{\sigma|0},\quad\det\sigma'=\det\sigma
\]
is conjugate \unsure{local\textsuperscript{ly}} to $\sigma$ --- and
\deleted{\unsure{it suffices for}} \inserted{the conditions}
are written (since
$\unsure{\det}(a\sigma+b)=a^2\det\sigma+
 ab\underbrace{\deleted{\unsure{\det}}\inserted{\operatorname{Tr}}\sigma}_{0}+b^2$
\[
\left\{\begin{aligned}
-va^2+b^2&=-v\\
2b&=0\qquad(\text{automatic in characteristic }2)
\end{aligned}\right.
\]
which in no way imply $b=0$ ....

\deleted{\illegible{1}} Let us now consider the product group
\[
B'_0\times A_\rho\times N'_\sigma,
\]

\scanpage{317}

and the morphism of schemes
\[
\tag{25}
\left\{\begin{aligned}
B'_0\times A_\rho\times N'_\sigma
&\longrightarrow \deleted{\mathcal G_{\rho,\sigma}}
 G\times G\times\unsure{\mathbb E^1|\mathbb E'}\times\unsure{\mathbb E^1|\mathbb E'}\\
(\underline u,\underline a,\underline b)
&\longmapsto
(\alpha=\underline u\underline a^{-1},\beta=\underline u\underline b^{-1},
\mu=\det\underline u,\varepsilon=\varepsilon_\sigma(\underline b)),
\end{aligned}\right.
\]

I say that the image lies in $\mathcal G^*_{\rho,\sigma}$, i.e.
that we will indeed have
\[
\tag{\unsure{*}}
\alpha\rho(\alpha)^{-1}
=\varepsilon\bigl(\beta\sigma(\beta)^{-1}\bigr)\varepsilon_1(\mu-1).
\]
Indeed
\[
\begin{aligned}
\alpha\rho(\alpha)^{-1}&=\underline u\rho(\underline u)^{-1}\\
\beta\rho(\beta)^{-1}
&=\underline u\underbrace{
\underline b^{\inserted{-1}}\sigma(\underline b^{-1})^{-1}
}_{\varepsilon_\sigma(\underline b)^{-1}=\varepsilon_\sigma(\underline b)}
\sigma(\underline u)^{-1}
=\varepsilon\underline u\sigma(\underline u)^{-1}
\end{aligned}
\]
and we know (by II) that
\[
\underline u\rho(\underline u)^{-1}
=\underline u\sigma(\underline u)^{-1}\varepsilon_1(\mu-1)
\]
hence $\unsure{(*)}$, since $\varepsilon^2=1$. But
moreover the morphism (25) is a monomorphism, \deleted{therefore}
for in terms of $(\alpha,\beta,\mu,\varepsilon)$, $u$ is determined
\inserted{(in terms of $\alpha$ already)},
(as the \uline{unique} section of $B'_0$ which
is of the form $\alpha\underline a^{-1}$, with $\underline a\in\Gamma T_\rho$
(for $T_\rho\cap B'_0=\{1\}$)), and $\underline b$ is then determined
by the relation $\beta=\underline u\underline b^{-1}$. We thus find
that (25) is an \uline{isomorphism}
\[
\tag{26}
\left\{\begin{aligned}
B'_0\times A_\rho\times N'_\sigma
&\overset{\sim}{\longrightarrow}\mathcal G^*_{\rho,\sigma}\\
(\underline u,\underline a,\underline b)
&\longmapsto
(\alpha=\underline u\underline a^{-1},\beta=\underline u\underline b^{-1},
\mu=\det\underline u,\varepsilon=\varepsilon_\sigma(\underline b))
\end{aligned}\right.
\]

\scanpage{318}
\mspage{459}

If we have two sections $(\underline u,\underline a,\underline b)$ and
$(\underline u',\underline a',\underline b')$ of
$B'_0\times T_\rho\times N'_\sigma$, under what
conditions do we have
\margin{\unsure{(27)}}
\[
\tag{27}
\widetilde\alpha=\widetilde\alpha',\quad
\widetilde\beta=\widetilde\beta'\ ?
\]
It clearly suffices to have
\[
\tag{28}
\underline a'=\lambda_1\underline a,\qquad
\underline b'=\lambda_2\underline b
\qquad \lambda_1,\lambda_2\in\Gamma G_{\mathrm m}
\]
and I say that these conditions are also necessary.
This comes from the fact that
$B'_0\hookrightarrow\unsure{PG}\simeq\unsure{GP(1)}$
is a monomorphism (let $\widetilde B'_0$ be the image of $B'_0$),
that \deleted{\unsure{the}} $A_\rho\to PG$ and
$N'_\sigma\longrightarrow PG$ have kernel $G_{\mathrm m}$
(\unsure{let|let} $\widetilde T_\rho$ and $\widetilde N'_\sigma$ be their images),
and that we have in $PG$
\[
\tag{29}
\widetilde T_\rho\cap\widetilde B'_0=\{1\},
\]
so that \deleted{\illegible{1}} $\widetilde{\underline u}$
\inserted{$\in\Gamma\widetilde B'_0$ and $\widetilde{\underline a}\in\Gamma\widetilde A_\rho$}
are defined \uline{uniquely} in terms of $\widetilde\alpha$
by the relation
\[
\deleted{\illegible{1}}\qquad
\widetilde\alpha=\widetilde{\underline u}\widetilde{\underline a}^{-1}
\]
-- and that then $\widetilde{\underline b}$ is defined
\inserted{in terms of $\widetilde\alpha,\widetilde\beta$}
by the relation
\[
\widetilde\beta=\widetilde{\underline u}\widetilde{\underline b}^{-1}.
\]

This shows that \deleted{\unsure{\(\mathcal G\)}} we have an \uline{isomorphism}
\[
\tag{30}
\left\{\begin{aligned}
\widetilde B'_0\times\widetilde A_\rho\times\widetilde N'_\sigma
&\overset{\sim}{\longrightarrow}
\widetilde{\mathcal G}^{*}_{\rho,\sigma}
\subset PG\times PG\times\unsure{\mathbb E^1|\mathbb E'}\\
(\widetilde{\underline u},\widetilde{\underline a},\widetilde{\underline b})
&\longmapsto
(\widetilde{\underline u}\widetilde{\underline a}^{-1},
\widetilde{\underline u}\widetilde{\underline b}^{-1},
\delta(\widetilde{\underline u}))
\end{aligned}\right.
\]

\scanpage{319}

where
\[
\tag{31}
\delta(\widetilde{\underline u})=\det(\underline u)
\]
when $\widetilde{\underline u}$ comes from $\underline u\in\Gamma B'_0$ --
or, what amounts to the same (cf. page 441)
$\delta(\widetilde{\underline u})$ is defined by
\[
\tag{32}
\widetilde{\underline u}(N_0)=\delta(\widetilde{\underline u})N_0.
\]

\margin{\uline{N.B.} The \unsure{centralizer} of $\widetilde\sigma$ in
$\widetilde G=PG$ \uline{is not} $\widetilde T_\sigma$, but
$\widetilde N'_\sigma$ \deleted{\illegible{2}}.}

\uline{Remark.} Let us note that a section $(\widetilde\alpha,\widetilde\beta,\mu)$
\deleted{\unsure{$\in\Gamma$}} of $\widetilde{\mathcal G}^{*}_{\rho,\sigma}$
\unsure{remains} \deleted{\unsure{entirely}} determined,
\inserted{\uline{modulo the action of} $\unsure{\widetilde N'_\sigma}$ \uline{on} $\beta$},
\unsure{to} \unsure{be} \unsure{precise}, in terms of $\widetilde\alpha$ --
\unsure{condition|conditions} to be satisfied by $\widetilde\alpha$, so that
there exist $\widetilde\beta,\mu$ \deleted{\illegible{1}} which
\unsure{complete} it to a section of $\widetilde{\mathcal G}^{*}_{\rho,\sigma}$,
being the following: the coefficient $D$ of
$\zeta=\alpha\rho(\alpha)^{-1}=\begin{pmatrix}A&B\\C&D\end{pmatrix}$
is invertible (which makes it possible to determine $\mu$ by $\mu=1/D$), then
we determine $\widetilde\eta$ \deleted{\illegible{1}} by
$\widetilde\eta=\widetilde\zeta\widetilde\varepsilon_1(\mu-1)$, and
finally we find \inserted{the} $\widetilde\beta$ \inserted{that correspond}
by the condition
\[
\widetilde\beta\widetilde\sigma(\widetilde\beta)^{-1}=\widetilde\eta.
\]

\margin{\uline{N.B.} $\widetilde\alpha$ depends on three parameters,
while $\widetilde{\mathcal G}^{*}_{\rho,\sigma}$
\uline{\unsure{is of|admits}} relative \unsure{dimension} $4=2+1+1$
(for $\widetilde{\mathcal G}^{*}_{\rho,\sigma}
\simeq\widetilde B'_0\times\widetilde T_\rho\times\widetilde N'_\sigma$).}

\margin{\uline{N.B.} We have lifted \unsure{locally$^{\mathrm t}$}
$\widetilde\alpha$ to $\alpha\in\Gamma G$,
$\zeta$ \unsure{does not} depend on the choice of lifting.}

Likewise, a section $(\alpha,\beta,\mu,\varepsilon)$
\inserted{of $\mathcal G^*_{\rho,\sigma}$}
is determined by \deleted{\unsure{$\alpha$ alone}}
\inserted{\deleted{\illegible{2}} $\alpha$ alone} modulo
\deleted{\illegible{2}} $N'_\sigma$ -- we first find
$\underline u,\underline a$ in terms

\scanpage{320}
\mspage{460}

of $\alpha$ as above, hence $\mu=\det\underline u$, hence
and it \deleted{suffices} \deleted{$\varepsilon\beta\sigma(\beta)^{-1}=\unsure{\alpha|\eta}$}
remains to determine $\beta$ as
$\beta=\underline u\underline b^{-1}$, with $\underline b\in\Gamma N'_\sigma$
arbitrary, $\varepsilon$ being defined then in terms of the choice of $\underline b$
as being $\varepsilon_\sigma(\underline b)$. The condition for
fixed $\unsure{\alpha_0|\alpha}$ to be completable by $(\beta,\mu,\varepsilon)$
to $(\alpha,\beta,\mu,\varepsilon)\in\Gamma\mathcal G^{\unsure{\alpha|*}}_{\rho,\sigma}$,
is always the condition that the coefficient $D$ \deleted{of $\zeta=\alpha\unsure{+}$}
\[
\tag{33}
\zeta=\alpha\rho(\alpha)^{-1}=\begin{pmatrix}A&B\\C&D\end{pmatrix}
\]
be invertible, which \inserted{makes it possible to} \unsure{will determine|determine} $\mu$ by
\[
\tag{34}
\mu=1/D.
\]

Let us put
\[
\tag{35}
\alpha=\begin{pmatrix}a&b\\c&d\end{pmatrix}
\]
and make $\alpha\rho(\alpha)^{-1}$ explicit; with
$\rho=\begin{pmatrix}0&1\\v&1\end{pmatrix}$

\margin{$\delta=\det\alpha=ad-bc$}
\[
\begin{aligned}
\alpha\rho(\alpha)^{-1}
&=\begin{pmatrix}a&b\\c&d\end{pmatrix}
\begin{pmatrix}0&1\\v&\deleted{\unsure{w}}\inserted{1}\end{pmatrix}
\frac1\delta\begin{pmatrix}d&-b\\-c&a\end{pmatrix}
\left(-\frac1v\right)
\begin{pmatrix}\deleted{\unsure{w}}\inserted{1}&-1\\-v&0\end{pmatrix}\\
&=-\frac1{\unsure{\delta}v}
\begin{pmatrix}bv&a+b\deleted{\unsure{w}}\\dv&c+d\deleted{\unsure{w}}\end{pmatrix}
\begin{pmatrix}d\deleted{\unsure{w}}+bv&-d\\-c\deleted{\unsure{w}}-av&c\end{pmatrix},
\end{aligned}
\]
hence
\[
\deleted{\unsure{=\frac1{\delta v}}}
\]
\[
\tag{36}
\left\{\begin{aligned}
(-\delta v)A&=\unsure{v}b(d+\unsure{vb|bv})\unsure{-}(a+b)(c+va)
&(-\delta v)B&=-vbd+c(a+b)\\
(-\delta v)C&=vd(d+vb)-(c+d)(c+va)
&(-\delta v)D&=-vd^2+cd+c^2,
\end{aligned}\right.
\]
the most essential quantity thus being
\[
\tag{37}
\left\{\begin{aligned}
(-\unsure{\delta}v)D&=-vd^2+cd+c^2=\det(d\rho+c),\\
\text{\deleted{\unsure{and after that|and after calculation}} where }&
\left\{\begin{aligned}
\deleted{\unsure{\alpha}}\ \delta&=\det\alpha=ad-bc,\\
v&\ \text{the \unsure{parameter|coefficient} which defines }
\rho=\begin{pmatrix}0&1\\v&1\end{pmatrix}
\end{aligned}\right.
\end{aligned}\right.
\]